% Quadratic forms over a base in Grothendieck's Géométrie élémentaire schématique (1992),
% BnF NAF 29093 (3). Rewritten on 8 November 2026 as a mathematics article on that volume alone
% (the earlier archival survey of the other boxes is in git history, commit cfa9d8f0).
% Rights: the BnF allows no transcription of the fonds online without the heirs' authorisation,
% and very short quotations. Nothing below reproduces the volume: its results are described in our
% own words and notation, and at most a few words are quoted, with their folio.
% NOT ON THE SITE until Michel approves (removed from the Method tab and docs/study/tei.mjs on
% 8 November 2026).
% Compile: tectonic -X compile article.tex
\documentclass[11pt,a4paper]{amsart}
\usepackage[T1]{fontenc}
\usepackage[utf8]{inputenc}
\usepackage{lmodern,microtype,amssymb,mathtools,booktabs,tabularx,array}
\usepackage[margin=2.7cm]{geometry}
\usepackage[hidelinks]{hyperref}

\newcommand{\q}[1]{\guillemotleft\,#1\,\guillemotright}
\newcommand{\cO}{\mathcal{O}}
\newcommand{\Sym}{\operatorname{Sym}}
\newcommand{\End}{\operatorname{End}}
\newcommand{\GL}{\mathrm{GL}}
\newcommand{\PGL}{\mathrm{PGL}}
\newcommand{\Ort}{\mathrm{O}}
\newcommand{\SO}{\mathrm{SO}}
\newcommand{\Isom}{\operatorname{\underline{Isom}}}
\newcommand{\Gm}{\mathbf{G}_{\mathrm m}}
\newcommand{\Pj}{\mathbf{P}}
\newcommand{\Z}{\mathbf{Z}}
\newcommand{\tr}{\operatorname{tr}}
\newcommand{\disc}{\operatorname{disc}}
\newcommand{\status}[1]{\par\smallskip\noindent\textit{In the volume.} #1\par\smallskip}
\newcolumntype{Y}{>{\raggedright\arraybackslash}X}
\setlength{\emergencystretch}{2em}

\theoremstyle{plain}
\newtheorem{theorem}{Theorem}[section]
\newtheorem{proposition}[theorem]{Proposition}
\theoremstyle{definition}
\newtheorem{definition}[theorem]{Definition}
\newtheorem{remark}[theorem]{Remark}
\newtheorem{question}{Question}

\title[Quadratic forms over a base in Grothendieck's 1992 volume]{Quadratic forms over a base in Grothendieck's\\ \emph{Géométrie élémentaire schématique} (1992)}
\author{Michel Hua}
\address{Commutator, Paris, France}
\email{michel@commutator.io}
\thanks{ORCID 0009-0002-8304-0683.}
\date{8 November 2026}
\subjclass[2020]{11E08, 11E81, 14L15; 01A60}
\keywords{quadratic forms over rings, half-discriminant, orthogonal group schemes, projective involutions, cross-ratio, EGA, Grothendieck}

\begin{document}
\begin{abstract}
Among the papers Grothendieck wrote at Lasserre, now at the Bibliothèque nationale de France (NAF~29093), box~3 holds a continuous mathematical text of 549 leaves dated from 19 August to 1 October 1992, the \emph{Géométrie élémentaire schématique}. It treats invertible modules, $\varepsilon$-symmetric and quadratic forms, discriminants, orthogonal structures, smooth relative quadrics, involutions and projective involutions, and quadratic groups over an arbitrary base scheme, with statements numbered as in EGA; on one leaf the author presents it as a chapter of EGA. Its subjects are those of two paragraphs of EGA~IV announced in 1965 and never published. This article sets out, section by section, the mathematics the volume covers, stated in modern terms over an arbitrary base and without dividing by~$2$, and marks for each statement what is known of it: the section titles and a few statements were read in the reading room, the body of the first fourteen sections was not. Where Grothendieck's own earlier treatment survives (an EGA~V typescript kept at the IHÉS and the Montpellier folder~77), his statements are given; where it does not, the classical statement the section title names is given with its source in print, and is not attributed to him. We list what the volume appears to add to the printed theory, as claims about the literature, and the questions that can only be settled at the BnF.
\end{abstract}
\maketitle

\section{Introduction}

\subsection{The volume}
The papers of Grothendieck's last years entered the Bibliothèque nationale de France in 2023 and are described as NAF~29093, \emph{Fonds Alexandre Grothendieck. \OE uvres.}, 1987--1999, forty-three boxes \cite{bnf-notice}. Almost all of it is filed under \emph{Réflexions sur la Vie et le Cosmos}. Box~3 is the exception: the \emph{Géométrie élémentaire schématique}, 581 leaves according to the notice, of which 549 are a continuous text in eighteen sections and 32 are undated drafts. The text is written in ink on both sides of the leaves, dated session by session from 19 August to 1 October 1992 (about a dozen leaves a day), with almost no crossings-out; its statements are numbered on three levels as in EGA (for instance Proposition~15.1.17, f.~168v, and Corollary~15.3.4, f.~175r). The volume does not mention physics. It stops inside \S\,18.8 (f.~517v); the eighteenth section, announced in the table with cross-ratios in its title, is not finished. On f.~378 the author describes what he is writing as \q{un (très gros) chapitre des EGA}; the Département des Manuscrits allowed this quotation.

The boxes are digitised but served only in the reading room, and photographs are not allowed. The library's rule, given on 29 September 2026, is that no transcription may be put online without the heirs' authorisation, and that very short quotations are allowed. Accordingly nothing here reproduces the volume. Its results are described in our own words and notation, and quotations are kept to a few words with their folio.

\subsection{Its place}
In a letter to Dieudonné of 29 September 1965, kept in the Montpellier fonds, Grothendieck gives the plan of the end of EGA~IV, with \q{24. Formes lisses, singularités quadratiques ordinaires} and \q{26. Résultant et discriminant} (folder~25, p.~134). Neither paragraph appeared; the printed EGA~IV stops at \S\,21. Two unpublished states of the subject survive before 1992. An EGA typescript kept at the IHÉS, undated (the IHÉS suggests 1969), whose heading is corrected by hand from IV to V, defines smooth forms and proves a smoothness criterion for quadratic forms over any base through a corrected discriminant \cite{ihes}. The Montpellier folder~77, begun in 1977, works with the same half-discriminant, proves its integrality by a graph expansion of the universal determinant (p.~68), and goes on to quadratic covers, affine quadrics and orthogonal groups over a base; folder~84 (around 1982--1986) treats binary forms and quadratic algebras over a base. The 1992 volume takes the subject up again from the beginning, in the language of EGA and without reference to any printed volume.

\subsection{What was read, and what this article can claim}\label{sub:evidence}
The volume was consulted in two sessions, on 23 and 29 October 2026. What was read is: the table of contents (three unnumbered leaves), in which the titles of sections 1--8, 10, 14 and 18 were noted; the second pochette (\S\S\,15--16) and the third (\S\S\,17--18), leafed through, with the numbering, a few dates and the notation noted; and the 25 leaves of drafts. The first pochette, which holds \S\S\,1--14 (ff.~1--155), was read for its table only. The titles and folio ranges of \S\S\,9--13, 15--16 and 17 below are those of the BnF's notice of the box \cite{bnf-notice}. No statement of \S\S\,1--14 was read on the page.

This settles what the article can do. It cannot reconstruct Grothendieck's proofs of 1992, nor say which of the classical results below he states, proves or merely uses. It can say precisely what mathematics each section covers, because the titles and notation fix it, and it can give his own earlier statements where a surviving text has them. Each subsection of Section~\ref{sec:math} therefore opens with what is known of the section in the volume (marked \emph{title}, \emph{notice} or \emph{seen}), followed by statements labelled by their source: a printed reference, the IHÉS typescript, folder~77, or the present author, in which case the statement is ours and is proved or checked here.

\subsection{Why it matters}
The layer the volume treats, quadratic forms on locally free modules over an arbitrary base, including the characteristic~$2$ fibres of odd rank, is now standard through Knus's book \cite{knus}, Conrad's appendix on orthogonal group schemes \cite{conrad}, and the study of quadratic and quaternion rings and their discriminants over an arbitrary base \cite{voight2011, voight2016, wood}. Conrad credits the half-discriminant, the device that makes the odd-rank theory work over $\Z$, to Grothendieck and Kneser independently \cite[\S\,C.1]{conrad}. The 1992 volume is the longest text by Grothendieck on this layer known to us, and more than half of it (\S\S\,15--18, ff.~156--549) is about involutions, projective involutions and quadratic groups, subjects that neither the EGA~IV plan nor the IHÉS typescript contains.

The other mathematics of the fonds read in October and November 2026 (a course on triality of 1995, notes on a combinatorial geography of 1996, and dated work on paths on a chessboard and a cube of 1997--1998) is left for another occasion.

\subsection*{Conventions}
$S$ is a scheme, $E$ a locally free $\cO_S$-module of finite rank $n$, $L$ an invertible $\cO_S$-module. A \emph{quadratic form} $q\colon E\to L$ is a map of sheaves with $q(cx)=c^2q(x)$ whose polar $b_q(x,y)=q(x+y)-q(x)-q(y)$ is bilinear; so $b_q(x,x)=2q(x)$, and nothing is divided by~$2$. Locally $q=\sum_{i\le j}a_{ij}x_ix_j$. The \emph{standard form} is $q_{2m}=x_1x_2+\dots+x_{2m-1}x_{2m}$ and $q_{2m+1}=x_0^2+q_{2m}$ \cite[Def.~C.1.2]{conrad}; the IHÉS typescript puts the square in the last variable, which changes nothing below.

\section{The mathematics, section by section}\label{sec:math}

\subsection{Invertible modules (\S\,1)}
\status{Title only.}
The section that opens the volume concerns invertible modules. The statement every later section needs is Kummer theory for $\mu_2$ in the fppf topology: the category of pairs $(N,\phi)$, $N$ invertible and $\phi\colon N^{\otimes2}\xrightarrow{\sim}\cO_S$, is equivalent to that of $\mu_2$-torsors, so that $H^1_{\mathrm{fppf}}(S,\mu_2)$ sits in $0\to\cO(S)^\times/\cO(S)^{\times2}\to H^1(S,\mu_2)\to\operatorname{Pic}(S)[2]\to0$. Square roots of a unit $\lambda$ form the torsor attached to $(\cO_S,\lambda^{-1})$. This torsor reappears below as the double cover of a quadratic form of odd rank (\S\,\ref{sub:covers}) and as the obstruction to lifting a projective involution (\S\,\ref{sub:projinv}).

\subsection{$\varepsilon$-symmetric bilinear forms and adjunction (\S\S\,2--3)}
\status{Titles only (\S\,2, $\varepsilon$-symmetric bilinear forms; \S\,3, adjunction).}
For $\varepsilon\in\mu_2(S)$, a bilinear $B\colon E\times E\to L$ is $\varepsilon$-symmetric if $B(y,x)=\varepsilon B(x,y)$. Over a base where $2$ is not invertible three notions separate: $(-1)$-symmetric, alternating ($B(x,x)=0$), and polar of a quadratic form (symmetric with $B(x,x)\in 2\cdot L$). The polar of $q$ is symmetric; in characteristic~$2$ it is alternating. When $B$ is perfect, each $u\in\End(E)$ has an adjoint $u^*$ with $B(ux,y)=B(x,u^*y)$, and $u\mapsto u^*$ is an anti-involution of $\End(E)$. Knus's book is written in the generality of hermitian forms over rings with involution, which contains these \cite{knus}.

\subsection{Modules of rank $2$, with and without involution (\S\S\,4--6)}
\status{Titles only: locally free modules of rank~$2$ (\S\,4), the same with an involution (\S\,5), and two modules $M^+_E$, $M^-_E$ regarded as quadratic modules (\S\,6).}
For $E$ of rank~$2$ the determinant is a perfect alternating pairing $E\times E\to\det E$; the adjoint it defines is the canonical involution $u\mapsto\bar u=\tr(u)-u$ of $\End(E)$, with $u\bar u=\det(u)$. Two quadratic modules come with it, and they are the natural candidates for the $M_E^\pm$ of \S\,6, which we have not seen defined:
\begin{enumerate}
\item $(\End(E),\det)$, of rank~$4$, locally the form $x_1x_4-x_2x_3$, hence smooth in every characteristic;
\item $(\End^0(E),\det)$, the trace-zero endomorphisms, of rank~$3$: on $\left(\begin{smallmatrix}x&y\\ z&-x\end{smallmatrix}\right)$ the determinant is $-(x^2+yz)$, a smooth ternary form, with half-discriminant a unit (compute with Theorem~\ref{thm:halfdisc}).
\end{enumerate}
Through (2), $\PGL(E)=\PGL_2$ acts on $\End^0(E)$ by conjugation preserving $\det$, which gives the isomorphism $\PGL_2\simeq\SO_3$ over $\Z$; Conrad gives it in this form \cite[Ex.~C.6.2]{conrad}. Through (1), $\GL(E)\times\GL(E)$ acts by $x\mapsto gxh^{-1}$, which gives $(\mathrm{SL}_2\times\mathrm{SL}_2)/\mu_2\simeq\SO_4$ \cite[Ex.~C.6.3]{conrad}. The same identification appears in Montpellier folder~82 ($\SO(3)\simeq$ the projective group of the line), dated around 1982.

\subsection{Quadrics as products of projective lines (\S\S\,7, 12)}
\status{Title of \S\,7 read with an uncertain word (an isomorphism between a product and a variety of quadrics); title of \S\,12 from the notice, ff.~109r--128r.}
The classical statement is the Segre isomorphism over an arbitrary base: for $E_1,E_2$ of rank~$2$, the map $\Pj(E_1)\times_S\Pj(E_2)\to\Pj(E_1\otimes E_2)$ identifies the product with the quadric of rank-one tensors, the zero scheme of the determinant on $\operatorname{Hom}(E_1^\vee,E_2)\simeq E_1\otimes E_2$, a form of type (1) above, valued in $\det E_1\otimes\det E_2$. Conversely, for a smooth quadric surface $X=V(q)\subset\Pj^3_S$ the two rulings are defined only over the finite étale double cover $Z_q$ given by the centre of the even Clifford algebra, and accordingly $\SO(q)/\mu_2$ is the Weil restriction from $Z_q$ to $S$ of a form of $\PGL_2$ \cite[Ex.~C.6.3]{conrad}. Which of these the volume proves is not known.

\subsection{The canonical quadratic form (\S\,8)}
\status{Title only.}
Not identified. Two classical forms deserve the name over any base: the determinant on $\End(E)$ for $E$ of rank~$2$, and, for $E$ of rank~$4$, the Pfaffian $\operatorname{Pf}\colon\Lambda^2E\to\Lambda^4E$, defined over $\Z$ by $\omega\wedge\omega=2\operatorname{Pf}(\omega)$ (on $\sum_{i<j}a_{ij}e_i\wedge e_j$ it is $a_{12}a_{34}-a_{13}a_{24}+a_{14}a_{23}$), a smooth form of rank~$6$ whose polar is the wedge pairing and whose quadric is the Grassmannian of planes in~$E$. The leaves of drafts are full of exterior products, which favours the second; this is a hypothesis.

\subsection{Discriminants (\S\,9)}\label{sub:disc}
\status{Notice: \S\,9, discriminants, ff.~36v--48v. Not read.}
The polar $b_q$ is a map $E\to E^\vee\otimes L$; its determinant is a section $\disc(q)$ of $(\det E)^{\otimes-2}\otimes L^{\otimes n}$, the \emph{discriminant}. In even rank it is the right invariant. In odd rank it vanishes on every fibre of characteristic~$2$, and the following replaces it.

\begin{theorem}[half-discriminant]\label{thm:halfdisc}
Let $n$ be odd and $Q=\sum_{i\le j}A_{ij}x_ix_j$ the universal form over $\Z[A_{ij}]$.
\begin{enumerate}
\item $\disc(Q)=2\,\delta'(Q)$ with $\delta'(Q)\in\Z[A_{ij}]$, and the content of $\delta'(Q)$ is~$1$.
\item Specialising, every $q$ on $E$ of rank~$n$ has a half-discriminant $\delta'(q)\in\Gamma\bigl((\det E)^{\otimes-2}\otimes L^{\otimes n}\bigr)$, compatible with base change, with $\delta'(q\circ g)=(\det g)^2\delta'(q)$ for $g\in\GL(E)$, and $\delta'(q\oplus q')=\delta'(q)\disc(q')$ when $q'$ has even rank.
\item $\delta'(q_{2m+1})=(-1)^m$.
\end{enumerate}
\end{theorem}

\noindent\emph{Sources.} In print: Conrad \cite[Prop.~C.1.4]{conrad}, by the defect of a form over a field of characteristic~$2$. In the IHÉS typescript: (1) is Lemma~16.5 (the content of $\disc(Q)$ is exactly $2$), and the section $\tilde d(q)$ of (2), which the typescript calls the \q{discriminant corrigé}, is Proposition~16.6 and Definition~16.7, with the product rule as Corollary~16.8, including the factor $2$ when both summands have odd rank \cite{ihes}. In folder~77, p.~68: $\disc(Q)$ is written as a sum over the graphs of permutations of $\{1,\dots,n\}$, and modulo~$2$ only the fixed-point-free involutions survive, which is impossible in odd size. The integrality has been checked in Lean for every odd $n$, by the route through alternating matrices \cite{hua2611}. For $n=3$ and $q=ax^2+by^2+cz^2+uyz+vxz+wxy$ one finds $\delta'(q)=4abc+uvw-au^2-bv^2-cw^2$, which is Voight's (half-)discriminant of a ternary form \cite[(3.2)]{voight2011}.

\begin{remark}[naming]\label{rem:naming}
The typescript suggests a change of terminology: it proposes to call the classical $\disc(q)$ the \emph{determinant} of $q$, and to keep \emph{discriminant} for $\tilde d$, on the ground that $\tilde d$ is what generalises to the discriminant of a form of any degree, the subject of the announced \S\,26 \cite{ihes}. The 1992 volume has a remark of its own on terminology, a point he says he had not thought of in the time of EGA and SGA (f.~156v); what that point is was not noted.
\end{remark}

\subsection{Orthogonal structures (\S\,10)}\label{sub:orth}
\status{Title read (\q{Digressions}, variations on orthogonal structures); notice: ff.~48v--69v. Not read.}
Call $q$ \emph{smooth} if for every $s\in S$ the quadric $V(q_s)\subset\Pj(E_s^\vee)$ is smooth of dimension $n-2$ and $q_s\neq0$. At a geometric point this means that the radical of $b_q$ is $0$, or is a line on which $q$ is non-zero (characteristic~$2$, $n$ odd). Conrad calls this non-degenerate, SGA~7 calls it ordinary, Knus calls it regular or semi-regular \cite[\S\,C.1 and Rem.~C.1.1]{conrad}; folder~77 calls it regular.

\begin{theorem}[orthogonal group schemes]\label{thm:orth}
Let $q\colon E\to\cO_S$ be smooth of rank $n$.
\begin{enumerate}
\item The functor $S'\mapsto\{g\in\GL(E_{S'}) : q\circ g=q\}$ is a closed subgroup scheme $\Ort(q)\subset\GL(E)$ of finite presentation. It is smooth if and only if $n$ is even or $2\in\cO_S^\times$, and then of relative dimension $n(n-1)/2$.
\item For $n$ odd, $\Ort(q)=\mu_2\times\SO(q)$ with $\SO(q)=\ker\det$, and $\SO(q)$ is smooth.
\item $\Isom(q_n,q)$ is an $\Ort_n$-torsor for the fppf topology, and $q\mapsto\Isom(q_n,q)$ is an equivalence between smooth forms of rank $n$ and $\Ort_n$-torsors. Hence isomorphism classes of pairs $(E,q)$, $E$ locally free of rank $n$ and $q$ smooth, are classified by $H^1_{\mathrm{fppf}}(S,\Ort_n)$.
\end{enumerate}
\end{theorem}

\noindent\emph{Sources.} (1) and (2): Conrad \cite[Thm.~C.1.5]{conrad}, through the orbit map $\GL(E)\to\{\text{smooth forms}\}$ and a tangent space computation. (3): stated in the IHÉS typescript (Remarques numbered 15.12, read as 16.12), the equivalence being said to follow from flat descent \cite{ihes}; the typescript also notes that $\Ort(q)$ is smooth near $s$ except exactly when $k(s)$ has characteristic~$2$ and $n$ is odd (16.13). The formalisation roadmap \cite{roadmap} sets out (1)--(3) on $\cO_S^n$ as targets.

\subsection{The double covers of a form (\S\,11)}\label{sub:covers}
\status{Notice: \S\,11, the coverings $\mathrm{Dis}(Q)$, ff.~69v--108v; the longest section of the first pochette. Not read.}
Two classical double covers are attached to a smooth form, one for each parity, and they are what a covering of $S$ named after the discriminant can be.
\begin{enumerate}
\item \emph{Odd rank $n=2m+1$.} By Theorem~\ref{thm:halfdisc}(2), a $g\in\Isom(q_n,q)$ satisfies $(\det g)^2\delta'(q)=(-1)^m$ on $\cO_S^n$. So $\det$ maps the $\Ort_n$-torsor of Theorem~\ref{thm:orth}(3) to the $\mu_2$-torsor of square roots of $(-1)^m\delta'(q)^{-1}$, a finite locally free double cover of $S$, étale only where $2$ is invertible. It is the obstruction to reducing the torsor from $\Ort_n=\mu_2\times\SO_n$ to $\SO_n$. The IHÉS typescript notes this reduction, with $\mu_2=\Ort(n)/\SO(n)$, as the way to obtain local isotriviality in the bad case (16.13) \cite{ihes}; folder~77 introduces the cover of square roots of $\delta'(q)$ (p.~105).
\item \emph{Even rank $n=2m$.} The centre $Z_q$ of the even Clifford algebra $C_0(q)$ is a finite étale $\cO_S$-algebra of rank~$2$, $C_0(q)$ is Azumaya over $Z_q$, and $\Ort(q)$ acts on $Z_q$ through its Galois group, which identifies $\Ort(q)/\SO(q)$ with $\Z/2$ \cite[Prop.~C.2.2 and \S\,C.2]{conrad}. When $2$ is invertible $Z_q$ is $\cO_S[t]/(t^2-(-1)^m\disc(q))$. Over $\Z$, $(-1)^m\disc$ of the universal form is congruent modulo~$4$ to the square of the hafnian of its off-diagonal coefficients, which yields an explicit free quadratic algebra $\cO_S[t]/(t^2-Bt-C)$ of discriminant $(-1)^m\disc(q)$ over any base. Folder~77 proves this congruence and uses it in this way; the roadmap \cite{roadmap} lists it as a target.
\end{enumerate}
Quadratic algebras over an arbitrary base, their discriminants and their relation to binary quadratic forms are the subject of Voight's and Wood's papers \cite{voight2016, wood}, after Bhargava's reformulation of Gauss composition over $\Z$ \cite{bhargava}; folder~84 treats the same correspondence between binary forms and quadratic algebras with an invertible module.

\subsection{Smooth relative quadrics (\S\,13)}\label{sub:smooth}
\status{Notice: \S\,13, smooth relative quadrics, ff.~128r--137v. Not read.}

\begin{theorem}[smoothness criterion]\label{thm:smooth}
For $q$ on $E$ of constant rank $n$ the following are equivalent: (i) $q$ is smooth; (ii) $\disc(q)$ (for $n$ even) or $\delta'(q)$ (for $n$ odd) is invertible; (iii) there is a surjective $S'\to S$ over which $q$ becomes isomorphic to $q_n$; (iii bis) the same with $S'\to S$ faithfully flat of finite presentation. If no residue field of characteristic~$2$ occurs together with odd rank, $S'\to S$ can be taken étale surjective, and finite étale locally on $S$.
\end{theorem}

\noindent\emph{Sources.} The IHÉS typescript states this as Theorem~16.9, with the étale refinement in 16.13, citing SGA~3, Exp.~XXIV, for the finite étale version \cite{ihes, sga3}. Its proof: $\Isom(q_n,q)$ is the fibre over $q$ of the orbit map $\GL(E)\to\mathcal Q(E)^*$ to the open scheme of forms with invertible (half-)discriminant; this map is faithfully flat by generic flatness and transitivity on geometric points; and over an algebraically closed field a form with non-zero half-discriminant is standard, by splitting off the radical line in characteristic~$2$. Conrad proves (i)$\Leftrightarrow$(ii) as part of Proposition~C.1.4(3), by the Witt-vector lift \cite{conrad}. The étale refinement fails in the excluded case. Over an imperfect field $k$ of characteristic~$2$ and $a\in k\smallsetminus k^2$, the form $at^2$ is smooth but becomes standard only over the inseparable extension $k(\sqrt a)$; the typescript gives this example (16.13).

The same typescript begins with smooth forms of any degree $d$: $\phi\in\Gamma(\Sym^dE)$ is smooth when $V(\phi)\subset\Pj(E)$ is smooth over $S$ and $\phi$ is fibrewise non-zero, which by the Jacobian criterion means that $\phi$ and its partial derivatives generate an ideal containing a power of the augmentation ideal. It also introduces the corresponding singularities, the \emph{elementary} singular points of multiplicity $d$ \cite{ihes}. Whether the 1992 volume returns to forms of degree $d>2$ was not seen (Question~\ref{qu:degree}).

\subsection{Orthogonal bifurcation (\S\,14)}
\status{Title only (ff.~137v--155, between the end of \S\,13 in the notice and the start of the second pochette).}
Not identified. We record the title and make no guess.

\subsection{Involutions and reflections (\S\,15)}\label{sub:refl}
\status{Seen (second pochette, ff.~156--245): statements numbered on three levels, for instance Proposition~15.1.17 (f.~168v) and Corollary~15.3.4 (f.~175r), dated 1 September 1992 at f.~190r. Notice: involutions, reflections. No statement read.}
The basic object over any base, in any characteristic, is the reflection. Let $q\colon E\to\cO_S$ and let $v\in E(S)$ with $q(v)\in\cO(S)^\times$. Put $s_v(x)=x-b_q(x,v)\,q(v)^{-1}v$. Then $q\circ s_v=q$, $s_v^2=1$ and $\det s_v=-1$.

\noindent\emph{Proof} (ours, elementary). With $t=b_q(x,v)/q(v)$, $q(x-tv)=q(x)-t\,b_q(x,v)+t^2q(v)=q(x)$. Next, $b_q(s_vx,v)=b_q(x,v)-2t\,q(v)=-b_q(x,v)$, so $s_v(s_vx)=x$. Finally $s_v=1-v\otimes\ell$ with $\ell(v)=2$, so $\det s_v=1-2$. In characteristic~$2$ the determinant is $1$ and $s_v$ is a transvection; there $\det$ does not detect $\Ort(q)/\SO(q)$ in even rank, and the Dickson invariant, defined through $Z_q$, takes its place \cite[(C.2.2)]{conrad}. Which properties of involutions in $\Ort(q)$ the section proves was not seen.

\subsection{Projective involutions (\S\S\,16--17)}\label{sub:projinv}
\status{Seen (\S\,16, ff.~246r--377v): projective bundles, the projective group, a set written $\mathrm{PsInv}(M)$, and at f.~297v a map from a set of involutions of a certain type (index not read) to the projective involutions, $\mathrm{Inv}_*(M)\to\mathrm{Invpr}(M)$. Notice: \S\,17, projective involutions, geometric study, ff.~378r--452v. The reading of $\mathrm{PsInv}$ below is a hypothesis.}
Let $M$ be locally free of rank $r$, $P=\Pj(M)$, so that $\operatorname{Aut}_S(P)=\PGL(M)$. A \emph{projective involution} is $\gamma\in\PGL(M)(S)$ with $\gamma^2=1$. The following is our reconstruction of what a map $\mathrm{Inv}_*(M)\to\mathrm{Invpr}(M)$ can measure; we do not attribute it to the volume.

\begin{proposition}\label{prop:lift}
Let $\gamma$ be a projective involution of $\Pj(M)$.
\begin{enumerate}
\item Zariski-locally on $S$, $\gamma$ lifts to $u\in\GL(M)$, and then $u^2=\lambda\cdot1$ with $\lambda\in\cO^\times$. The lifts form a $\Gm$-torsor $T_\gamma$, and $u\mapsto u^2$ satisfies $(cu)^2=c^2u^2$.
\item The lifts of $\gamma$ that are involutions of $M$ ($u^2=1$) form a $\mu_2$-torsor $I_\gamma$ for the fppf topology, namely the one attached by Kummer theory to $T_\gamma$ and the squaring map. Thus $\gamma$ lifts to an involution of $M$ over $S$ if and only if the class of $I_\gamma$ in $H^1_{\mathrm{fppf}}(S,\mu_2)$ vanishes, and two such lifts differ by $\pm1$.
\end{enumerate}
\end{proposition}

\begin{proof}
The sequence $1\to\Gm\to\GL(M)\to\PGL(M)\to1$ is exact for the Zariski topology on $M$ free, and $\gamma^2=1$ gives $u^2\in\Gm$. If $u$ and $u'=cu$ both square to $1$, then $c^2=1$. A unit $\lambda$ is fppf-locally a square, so lifts with square $1$ exist fppf-locally. The torsor of square roots of $\lambda$ along the transition functions of $T_\gamma$ is $I_\gamma$.
\end{proof}

Elements $u$ with $u^2$ a unit scalar, taken modulo $\Gm$, are thus the same as projective involutions, locally. This is our reading of $\mathrm{PsInv}(M)$. When $2$ is invertible and $u^2=1$, $M=M^+\oplus M^-$, and the fixed scheme of $\gamma$ is $\Pj(M^+)\sqcup\Pj(M^-)$. In characteristic~$2$, a lift with $u^2=1$ has $(u-1)^2=0$: it is unipotent and its fixed scheme is the projective bundle of $\ker(u-1)$ when that is a direct summand. A \emph{geometric study} in \S\,17 would be the study of these fixed loci and of their ranks over a base where both cases occur. Over fields, the involutions that Knus, Merkurjev, Rost and Tignol study are anti-automorphisms of central simple algebras \cite{kmrt}, a different notion from these elements of order~$2$ of $\PGL(M)$.

\subsection{Quadratic groups and cross-ratios (\S\,18)}
\status{Table: \S\,18 has \q{relations anharmoniques} (cross-ratios) in its title. Seen (ff.~453r--549r): finite monoids, monogenic groups, a notation $\mathrm{Scind}_S(\varepsilon)$ (perhaps the splittings of $\varepsilon$ over $S$), exterior products; \S\,18.8, on dualities, begins at f.~517v (18.8.2 at f.~532v), and the text stops inside it. Drafts (fourth pochette): references to \S\S\,18.7--18.11, one of them on divisors with cross-ratio relations, and a leaf on cross-ratio and duality computing in the algebra of matrices $\left(\begin{smallmatrix}a&b\\ c&d\end{smallmatrix}\right)$ with $a,d\in\Gamma(\cO_S)$, $b\in\Lambda$, $c\in\Lambda^{-1}$ for an invertible $\Lambda$; computations only, no numbered statement.}
That algebra is $\End(\cO_S\oplus\Lambda)$, whose units act on the $\Pj^1$-bundle $\Pj(\cO_S\oplus\Lambda)$. So the planned end of the volume is the projective line over a base, with its group and its cross-ratio. The classical statement is the following. Three pairwise disjoint sections of a $\Pj^1$-bundle $P\to S$ (disjoint on every fibre) trivialise it, $P\simeq\Pj^1_S$ with the sections at $0,1,\infty$, uniquely. A fourth section disjoint from the three is then a unit $\lambda$ with $1-\lambda$ a unit, the cross-ratio, that is, a morphism $S\to\Pj^1\smallsetminus\{0,1,\infty\}$; and $\mathfrak S_4$ acts on the cross-ratio through $\mathfrak S_3$, by $\lambda\mapsto1-\lambda,1/\lambda,\dots$. Which divisors with cross-ratio relations he had in mind is not known; the sections that would say were not written. What \emph{quadratic groups} denotes is not settled either: orthogonal groups, or group schemes of rank~$2$ over $S$ such as $\mu_2$, $\Z/2$ and their twisted forms, whose classification over a base goes back to Tate and Oort \cite{tateoort}. The vocabulary of finite monoids and monogenic groups points to the second reading (Question~\ref{qu:quadgroups}).

\section{What is new or different compared with print}\label{sec:new}

The claims below are about the printed literature consulted, not about priority.

\begin{enumerate}
\item \emph{The text itself.} We know of no printed or online description of the mathematics of the 1992 volume beyond the BnF's notice and this project's own notes. In particular the attribution of the half-discriminant to Grothendieck in \cite[\S\,C.1]{conrad} rests, in print, on no text of his. Two unpublished texts of his on this layer can be read today, the IHÉS typescript and folder~77, and the 1992 volume is a third and much longer one.
\item \emph{The shape of the theory.} In Conrad's appendix and in the roadmap \cite{roadmap}, the theory runs from forms to discriminants to orthogonal groups to Clifford algebras. The 1992 volume, judging by its table, starts lower: invertible modules, $\varepsilon$-symmetric forms valued in invertible modules, adjunction, and modules of rank~$2$ with involution all come before the discriminant (\S\S\,1--8, ff.~1--36). It also gives more than two thirds of its length to involutions, projective involutions and quadratic groups (\S\S\,15--18, about 390 of 549 leaves). No section on Clifford algebras appears among the titles known. Nothing in Conrad's Appendix~C, in the roadmap, or in Knus's book as far as we consulted it treats projective involutions of $\Pj(M)$ over a base as a subject (Proposition~\ref{prop:lift} is ours).
\item \emph{Terminology.} The IHÉS typescript's proposal to rename the classical discriminant \q{déterminant} and keep \q{discriminant} for $\tilde d$ (Remark~\ref{rem:naming}) has not been adopted: Conrad keeps the discriminant for the determinant and calls $\tilde d$ the half-discriminant. Whether the 1992 volume keeps the proposal is part of Question~\ref{qu:disc}.
\item \emph{The relative cross-ratio.} The planned end of the volume, cross-ratios over a base computed in $\End(\cO_S\oplus\Lambda)$, ties the quadratic theory to the projective line over a base, the object of $M_{0,4}$ and of his work of the 1980s on moduli of curves. The volume does not make this link in the leaves read; we note it as a direction the table points to, not as a result.
\end{enumerate}

\section{What remains open or unreadable}\label{sec:open}

The first fourteen sections, the half that overlaps the IHÉS typescript, were not read beyond their titles, and the titles of \S\S\,9, 11, 12, 13, 15, 16 and 17 are known only from the notice. Two titles were read with an uncertain word (\S\,7, and in \S\,10 the word read as \q{moteur}). The last section is unfinished, and \S\S\,18.9--18.11 exist only as references in the drafts. Nothing in the volume can be checked from outside the reading room. The questions below are those a further session would have to answer, in this order.

\begin{question}\label{qu:disc}
\S\,9 (ff.~36v--48v) and \S\,11 (ff.~69v--108v): how is the discriminant normalised in odd rank? Is it the half-discriminant, is it proved integral, and by which argument (content of the universal determinant as in the typescript, a graph expansion as in folder~77, or a defect argument)? Where does it take its values, $(\det E)^{\otimes-2}\otimes L^{\otimes n}$ or some other module? Does the volume keep the typescript's \q{déterminant}?
\end{question}

\begin{question}
\S\,13 (ff.~128r--137v): is there a theorem of the form of Theorem~\ref{thm:smooth}, with which topology for local triviality (fppf, étale, finite étale), and with the imperfect-field example?
\end{question}

\begin{question}
\S\,10 (ff.~48v--69v): are $\Ort(q)$ and $\SO(q)$ defined as group schemes, are smoothness and the torsor classification stated, and how is $\SO$ defined in even rank over a base where $2$ is not invertible (Dickson invariant, Clifford centre, or something else)?
\end{question}

\begin{question}
\S\S\,4--8 (ff.~1--36): what are $M^+_E$ and $M^-_E$? Which isomorphism between a product and a variety of quadrics does \S\,7 state? Which form is \emph{the} canonical quadratic form of \S\,8?
\end{question}

\begin{question}\label{qu:degree}
Does the volume treat forms of degree $d>2$, the resultant, or the discriminant of a form of degree $d$ (the subject of the announced EGA~IV \S\,26)?
\end{question}

\begin{question}
\S\,14 (ff.~137v--155): what is \emph{bifurcation orthogonale}?
\end{question}

\begin{question}
\S\S\,15--17: what are $\mathrm{PsInv}(M)$ and $\mathrm{Inv}_*(M)$ exactly, what is the index $*$ at f.~297v, and does the map to $\mathrm{Invpr}(M)$ come with an obstruction in $H^1(S,\mu_2)$ as in Proposition~\ref{prop:lift}? What are the main statements around Proposition~15.1.17 (f.~168v)?
\end{question}

\begin{question}\label{qu:quadgroups}
\S\,18: what is a \emph{quadratic group}, and what is $\mathrm{Scind}_S(\varepsilon)$? What does \S\,18.8 dualise?
\end{question}

\begin{question}
f.~156v: which point of terminology does he regret not having seen in the time of EGA and SGA? f.~378: how many years does he say he had left the chapter aside?
\end{question}

The volume is a clean text with numbered statements, and a reader trained in quadratic forms over a base could read its first pochette in one or two sessions. The heirs' authorisation would be needed to quote its statements.

\subsection*{Use of language models}
The questions for the sessions at the BnF were prepared, and the notes taken in the reading room organised, with language models (Claude Opus~5.5, Anthropic). The volume was read by the author, in person, and the readings reported here are his. The models also compared the notes with the IHÉS typescript, folder~77, Conrad's appendix and the formalisation roadmap, wrote Proposition~\ref{prop:lift} and the elementary computations of Section~\ref{sec:math} that carry no other source, checked the references against the publishers' records, and produced the text under the author's direction.

\begin{thebibliography}{99}
\bibitem{bhargava} M.~Bhargava, Higher composition laws~I: A new view on Gauss composition, and quadratic generalizations, \emph{Ann.\ of Math.~(2)} 159 (2004), 217--250, \href{https://doi.org/10.4007/annals.2004.159.217}{doi:10.4007/annals.2004.159.217}.
\bibitem{bnf-notice} Bibliothèque nationale de France, Département des Manuscrits, NAF~29093, \emph{Fonds Alexandre Grothendieck. \OE uvres.}, 1987--1999, 43 boxes; notice of box~3, \url{https://archivesetmanuscrits.bnf.fr/ark:/12148/cc126317b/cb5}, read in October 2026.
\bibitem{conrad} B.~Conrad, Reductive group schemes, in \emph{Autour des schémas en groupes~I}, Panoramas et Synthèses 42--43, Société mathématique de France, 2014; Appendix~C, Orthogonal group schemes. Numbering as in the author's version, \url{https://math.stanford.edu/~conrad/papers/luminysga3.pdf}.
\bibitem{hua2611} M.~Hua, Grothendieck Montpellier. Énoncés des lectures modernisées démontrés, Hexagon preprint \href{https://hexagonmath.org/2611.00125v1}{hexagon:2611.00125v1}, 2026; Lean file \texttt{lean/Grothendieck/Folder77.lean} of \url{https://github.com/Commutator-IO/grothendieck-project}.
\bibitem{ihes} A.~Grothendieck, typescript for EGA (heading corrected from IV to V), \q{Formes lisses} and \q{Formes quadratiques lisses}, undated (\q{1969 ?} in the IHÉS record), 18 leaves; Archives de l'IHÉS, fonds Grothendieck--Illusie, FG\_IL\_4\_EGA5\_3\_3.
\bibitem{kmrt} M.-A.~Knus, A.~Merkurjev, M.~Rost, J.-P.~Tignol, \emph{The Book of Involutions}, Amer.\ Math.\ Soc.\ Colloquium Publ.~44, 1998, \href{https://doi.org/10.1090/coll/044}{doi:10.1090/coll/044}.
\bibitem{knus} M.-A.~Knus, \emph{Quadratic and Hermitian Forms over Rings}, Grundlehren der mathematischen Wissenschaften~294, Springer, 1991, \href{https://doi.org/10.1007/978-3-642-75401-2}{doi:10.1007/978-3-642-75401-2}.
\bibitem{roadmap} M.~Hua, Roadmap: quadratic forms over a base, pull request TauCetiProject/TauCetiRoadmap\#1930 (open, November 2026), \url{https://github.com/TauCetiProject/TauCetiRoadmap/pull/1930}.
\bibitem{sga3} M.~Demazure, A.~Grothendieck (eds.), \emph{Schémas en groupes~III: Structure des schémas en groupes réductifs} (SGA~3, exposés XIX--XXVI), Lecture Notes in Math.~153, Springer, 1970, \href{https://doi.org/10.1007/BFb0059027}{doi:10.1007/BFb0059027}.
\bibitem{tateoort} J.~Tate, F.~Oort, Group schemes of prime order, \emph{Ann.\ Sci.\ École Norm.\ Sup.~(4)} 3 (1970), 1--21, \href{https://doi.org/10.24033/asens.1186}{doi:10.24033/asens.1186}.
\bibitem{voight2011} J.~Voight, Characterizing quaternion rings over an arbitrary base, \emph{J.\ reine angew.\ Math.} 657 (2011), \href{https://doi.org/10.1515/crelle.2011.054}{doi:10.1515/crelle.2011.054}, \href{https://arxiv.org/abs/0904.4310}{arXiv:0904.4310} (equation numbering of the arXiv version).
\bibitem{voight2016} J.~Voight, Discriminants and the monoid of quadratic rings, \emph{Pacific J.\ Math.} 283 (2016), 483--510, \href{https://doi.org/10.2140/pjm.2016.283.483}{doi:10.2140/pjm.2016.283.483}.
\bibitem{wood} M.~M.~Wood, Gauss composition over an arbitrary base, \emph{Adv.\ Math.} 226 (2011), 1756--1771, \href{https://doi.org/10.1016/j.aim.2010.08.018}{doi:10.1016/j.aim.2010.08.018}.
\end{thebibliography}
\end{document}
